\section{Conclusion and Future Work}
\label{sec:conclusion}

In this paper, we introduced \textbf{Q-PolyMerge}, a parameter-efficient, robust framework for quantization-aware test-time model merging on edge devices. By projecting layer-wise merging coefficients onto a low-degree continuous polynomial subspace of normalized layer depth, Q-PolyMerge resolves the \emph{Overfitting-Optimizer Paradox} that plagues traditional unconstrained test-time adaptation methods. 

Our approach reduces the learnable parameter space by over \textbf{78\%}, filtering out high-frequency optimization noise and mitigating the risk of degenerate entropy-collapse states. This continuous polynomial constraint not only stabilizes first-order gradient descent via the Straight-Through Estimator (STE) but also enables zero-order derivative-free search (1+1 Evolution Strategy) to run exceptionally well under standard INT8 regimes. Because 1+1 ES bypasses backpropagation and activation caching, it offers a zero-gradient-memory footprint, making test-time adaptation physically viable on low-power microcontrollers for the first time.

Extensive evaluation on a Vision Transformer backbone across four classification benchmarks under 8-bit and 4-bit post-training quantization demonstrates the superior robustness of Q-PolyMerge. Under 4-bit PTQ, Q-PolyMerge (Adam STE) achieves an average accuracy of \textbf{48.87\%}, strictly outperforming unconstrained Q-Merge (\textbf{46.02\%}) and naive post-merge quantization (\textbf{42.92\%}). Under 8-bit PTQ, Q-PolyMerge (Adam STE) achieves \textbf{59.76\%} average accuracy, which is highly comparable to the unconstrained Q-Merge baseline (\textbf{60.03\%}) while reducing standard deviation by over 47\% and recovering the unquantized continuous PolyMerge ceiling (\textbf{61.00\%}). Qualitative analysis of learned coefficient profiles confirms that Q-PolyMerge learns smooth, stable, and physically meaningful trajectories compared to the highly noisy and jagged trajectories of unconstrained methods.

\textbf{Future Directions and Scaling:} Scaling Q-PolyMerge to multi-billion parameter foundation models (e.g., LLaMA or CLIP) and deploying our zero-order pathway on physical hardware-in-the-loop testbeds (such as ARM Cortex-M7 or RISC-V GAP8) represent exciting avenues of research. Due to space constraints, we provide an in-depth discussion of these future directions, including depth scalability via B-splines and hardware deployment strategies, in Appendix~\ref{sec:appendix_future_work}.